\section{Graph Storage and Indexing}
\label{sec:gr-storage}

Once triples are extracted and normalized they must be stored in a system that
supports efficient traversal and querying, and the choice of data model shapes
what retrieval can do~\cite{angles2008graphdb}. Graph database models have been
surveyed extensively, and they differ in how they represent nodes, edges, and the
attributes attached to them~\cite{angles2008graphdb}. For Graph RAG the
requirement is twofold: the store must both hold the relational structure and
expose it to the retriever at query time~\cite{peng2024graphrag}.

\subsection{Triple Stores vs. Property Graphs}

Two dominant paradigms exist. RDF triple stores represent data strictly as
(subject, predicate, object) statements aligned with the W3C
standards~\cite{w3c2014rdf}, which maximizes interoperability and standardized
reasoning. Property graphs instead attach arbitrary key--value properties to both
nodes and edges, a model popularized for its flexibility and developer
ergonomics in connected-data applications~\cite{robinson2015graphdb}. Surveys of
graph data models frame these as points on a spectrum of expressiveness and
standardization rather than mutually exclusive
choices~\cite{angles2008graphdb, hogan2021kg}.

\subsection{Hybrid Vector-Graph Indexes}

Modern Graph RAG systems index the graph alongside the vector embeddings used by
dense retrieval, so that a query can exploit both semantic similarity and
explicit structure~\cite{peng2024graphrag}. LightRAG, for example, incorporates
graph structures into the text index and combines them with vector
representations to retrieve related entities and their relationships efficiently,
using an incremental update algorithm to keep the index current as new data
arrives~\cite{guo2024lightrag}. Such hybrid indexes are what let a single system
serve both the local, similarity-driven and the global, structure-driven retrieval
patterns~\cite{peng2024graphrag, guo2024lightrag}.
